\section{An attained envelope for mixed products}\label{sec:mixed}

Let $(\sigma_c)_c$ be independent uniform permutations of $[n]$, with $n\ge2$, and put $M_{cij}=\1_{\{\sigma_c(i)=j\}}$. A \emph{cell} is a triple $(c,i,j)$. For a nonempty set $E$ of distinct cells, assign integers $p_e,q_e\ge0$ with $p_e+q_e\ge1$, and set
\begin{equation}\label{eq:product}
P=\prod_{e\in E}(M_e-n^{-1})^{p_e}M_e^{q_e},\qquad a=|E|\le n.
\end{equation}
A cell is \emph{consistent} if no other cell of its color in $E$ shares its row or column. A \emph{centered singleton} has $(p_e,q_e)=(1,0)$. Let $h$ and $b$ count the consistent and inconsistent centered singletons, respectively. Consistency refers to the whole requested set, before a permutation is sampled.

Write $(n)_r=n(n-1)\cdots(n-r+1)$, with $(n)_0=1$, and define
\begin{equation}\label{eq:depletedpoly}
\mathcal P_j(n,r)
=j![z^j]\exp\left(\sum_{\ell\ge1}
\frac{(\ell-1)n+r}{\ell}z^\ell\right).
\end{equation}
Only a finite coefficient is extracted. The first polynomials are
$\mathcal P_0=1$, $\mathcal P_1=r$ and $\mathcal P_2=r^2+r+n$.

\begin{theorem}[An attained mixed-product envelope]\label{thm:mixed}
For the product \eqref{eq:product},
\begin{equation}\label{eq:mixedexact}
|\E P|\le 2^b\,\frac{n^{-h}\mathcal P_h(n,a-h)}{(n)_a}.
\end{equation}
For each fixed $n,a,h$, the right side with $b=0$ is the exact maximum over all permitted patterns, colors and multiplicities having these counts and no inconsistent centered singleton. It is attained by $a$ compatible cells in one color, with $h$ centered singletons and $a-h$ raw factors.
\end{theorem}

The multiplicities have no upper bound. At $a=1$, the upper bound is zero for a centered singleton and $1/n$ for every other allowed product. The scalar consequences, including the optimal $\sqrt e$ prefactor, are stated in Corollary~\ref{cor:mixedscalar}. When $n=1$ and $a\le n$, centered factors vanish and a purely raw product equals one.

\subsection{Exact depleted moments}

Fix one color and prescribe $r$ \emph{compatible} matches, meaning that their rows and columns are distinct. Their probability is $1/(n)_r$. Deleting those rows and columns leaves a uniform permutation matrix $U$ on $N=n-r$ symbols. Retain the original centering by putting $Z_{ij}=nU_{ij}-1$, whose row sum is $r$, and define
\[
\mu_j(n,r)=\E\prod_{i=1}^j Z_{ii},\qquad 0\le j\le N,\qquad \mu_0=1.
\]
Any remaining $j$ cells with distinct rows and columns have this moment after relabeling.

\begin{lemma}[Depleted row-sum recurrence]\label{lem:depleted}
For $1\le j\le n-r$,
\begin{equation}\label{eq:depleted}
(n-r-j+1)\mu_j=(r+2j-2)\mu_{j-1}+(j-1)\mu_{j-2},
\end{equation}
where the last term is omitted for $j=1$. Every $\mu_j$ is nonnegative.
\end{lemma}
\begin{proof}
Multiply the last row sum $\sum_v Z_{jv}=r$ by the preceding $j-1$ diagonal factors and take expectations. Each of the $N-j+1$ unused columns contributes $\mu_j$. For a used column $i<j$, the conflicting indicators have product zero, so
\[
Z_{ii}Z_{ji}=-Z_{ii}-Z_{ji}-1.
\]
After multiplying by the other diagonal factors, the two single-factor terms each contribute $-\mu_{j-1}$, by row relabeling, and the last contributes $-\mu_{j-2}$. Summing proves \eqref{eq:depleted}. Its denominator is positive. Induction starts with $\mu_0=1$ and, when $r<n$, $\mu_1=r/(n-r)$. At full depletion $r=n$, only the empty moment $\mu_0$ is used.
\end{proof}

For $r=0$, this is the recurrence of Collins and Nagatsu after rescaling \cite[Proposition 3.3, equation (16)]{CN}. To retain the exact moment, expand its indicators:
\[
\mu_j(n,r)=\sum_{\ell=0}^j
\binom j\ell(-1)^{j-\ell}\frac{n^\ell}{(n-r)_\ell}.
\]
Multiplication by $(n-r)_j/j!$ gives the coefficient identity
\begin{align*}
\frac{(n-r)_j\mu_j(n,r)}{j!}
&=\sum_{\ell=0}^j\frac{n^\ell}{\ell!}
(-1)^{j-\ell}\binom{n-r-\ell}{j-\ell}\\
&=[z^j](1-z)^{n-r}\exp\left(\frac{nz}{1-z}\right).
\end{align*}
Expanding the logarithm of the last generating function proves
\begin{equation}\label{eq:poly-moment}
\mu_j(n,r)=\frac{\mathcal P_j(n,r)}{(n-r)_j},
\qquad j+r\le n.
\end{equation}
This is also the classical Laguerre identity
$\mathcal P_j(n,r)=j!L_j^{(-n+r-1)}(-n)$ \cite[equation 18.12.13]{DLMF}. The coefficient derivation makes the following argument independent of that identification. Coefficients outside the moment range are only formal polynomials.

Let $\Pi_j$ denote the set partitions of $[j]$. The exponential formula in \eqref{eq:depletedpoly} reads
\begin{equation}\label{eq:depletedpartitions}
\mathcal P_j(n,r)=
\sum_{\pi\in\Pi_j}\prod_{B\in\pi}
\bigl((|B|-1)n+r\bigr)(|B|-1)!.
\end{equation}
Indeed, partitions with $m_\ell$ blocks of size $\ell$ number
$j!/\prod_\ell(\ell!)^{m_\ell}m_\ell!$. Multiplying their block weights gives exactly the coefficient in \eqref{eq:depletedpoly}.

\subsection{Merging colors and handling repeated cells}

Put
\[
F_n(r)=\frac{n^r}{(n)_r},\qquad
K_{j,r}=F_n(r)\mu_j(n,r)
=\frac{n^r\mathcal P_j(n,r)}{(n)_{j+r}}.
\]
\begin{lemma}[Monotonicity and merging]\label{lem:mixedmerge}
For fixed $n,j$, both $\mu_j(n,r)$ and $K_{j,r}$ are nondecreasing in $r$ on $0\le r\le n-j$. If $j_1+r_1+j_2+r_2\le n$, then
\[
K_{j_1,r_1}K_{j_2,r_2}\le K_{j_1+j_2,r_1+r_2}.
\]
\end{lemma}
\begin{proof}
Each block weight in \eqref{eq:depletedpartitions} increases with $r$, while $(n-r)_j$ decreases. Also the successive ratio of $n^r/(n)_{j+r}$ is $n/(n-j-r)\ge1$ whenever the next index is permitted.

For merging, unite partitions of two disjoint labeled sets and increase every block's parameter to $r_1+r_2$. Positivity gives
$\mathcal P_{j_1}(n,r_1)\mathcal P_{j_2}(n,r_2)
\le\mathcal P_{j_1+j_2}(n,r_1+r_2)$.
Finally, $(n)_u(n)_v\ge(n)_{u+v}$ when $u+v\le n$. These two comparisons prove the kernel inequality.
\end{proof}

\begin{proof}[Proof of Theorem~\ref{thm:mixed}]
Let $T$ be the $h$ consistent centered singletons. For each $e\notin T$, use $M_e^2=M_e$ to write its factor as $c_e+d_eM_e$. If $q_e\ge1$, then $c_e=0$ and $d_e=(1-1/n)^{p_e}\le1$. If $q_e=0$, then
\[
c_e=(-1/n)^{p_e},\qquad
d_e=(1-1/n)^{p_e}-(-1/n)^{p_e}\ge0.
\]
The coefficient weights satisfy
\begin{equation}\label{eq:weights}
d_e+n|c_e|\le
\begin{cases}
2,&(p_e,q_e)=(1,0),\\
1,&\text{otherwise}.
\end{cases}
\end{equation}
For the repeated centered case, put $x=1/n\le1/2$. At even $p\ge2$, the weight is at most
$(1-x)^2+(n-1)x^2=1-x$. At odd $p\ge3$, it is at most
$(1-x)^3+(n+1)x^3=1-3x+4x^2\le1-x$.
Thus repeated multiplicities cost no additional factor.

Expand the affine factors by choosing a set $S\subseteq E\setminus T$ of indicators. Incompatible choices vanish. For a compatible choice, let $r_c$ count its cells and $h_c$ the cells of $T$ in color $c$. Consistency keeps every cell of $T$ on the unexposed block. Exact conditioning gives
\begin{equation}\label{eq:prescribed}
\E\left[\prod_{e\in S}M_e
\prod_{e\in T}(M_e-n^{-1})\right]
=n^{-h-|S|}\prod_cK_{h_c,r_c}\ge0.
\end{equation}
The total index in each merge is at most $h+|S|\le a\le n$. Lemma~\ref{lem:mixedmerge} therefore gives
\[
\prod_cK_{h_c,r_c}\le K_{h,|S|}\le K_{h,a-h}.
\]
Taking absolute values of the affine coefficients and summing all subsets yields
\[
|\E P|\le n^{-a}K_{h,a-h}
\prod_{e\notin T}(d_e+n|c_e|)
\le 2^b n^{-a}K_{h,a-h}.
\]
Equation \eqref{eq:poly-moment} identifies this with \eqref{eq:mixedexact}.

For attainment, use the cells $(i,i)$, $1\le i\le a$, in one color. Center the first $h$ factors and take the other $r=a-h$ factors raw. Conditioning on the raw matches gives exactly
\[
\E P=\frac{n^{-h}\mu_h(n,r)}{(n)_r}
=\frac{n^{-h}\mathcal P_h(n,r)}{(n)_a}.
\]
This includes $h=0$, with only the empty moment at full depletion, and $a=h=1$, when the expectation is zero.
\end{proof}

\subsection{Scalar consequences and the optimal prefactor}

For $1\le h\le a\le n$, define
\begin{equation}\label{eq:mixedgamma-def}
\Gamma_{n,a,h}=
\frac{a+h-2+\sqrt{(a+h-2)^2+4(h-1)(n-a+1)}}
{2(n-a+1)}.
\end{equation}
\begin{lemma}[A geometric supersolution]\label{lem:super}
If $r+h\le a\le n$ and $h\ge1$, then
$0\le\mu_h(n,r)\le\Gamma_{n,a,h}^{\,h}$.
For $a\ge2$, put
\[
\beta=\frac{\sqrt{a-1}}{\sqrt n-\sqrt{a-1}}.
\]
Then $\Gamma_{n,a,h}\le\beta$, strictly if $h<a$, and equality holds if $h=a$. At $h=1$, $\Gamma_{n,a,1}=(a-1)/(n-a+1)$ is the exact moment at depletion $a-1$. For $h=0$, use $\mu_0=1$.
\end{lemma}
\begin{proof}
By monotonicity it suffices to use $r=a-h$. Put
$d=n-a+1$, $A=a+h-2$ and $v=h-1$. The number $\Gamma=\Gamma_{n,a,h}$ satisfies
$d\Gamma^2=A\Gamma+v$. If $\Gamma=0$, then $a=h=1$ and $\mu_1(n,0)=0$. Otherwise,
$\Gamma\ge A/d\ge r/(n-r)$.
For $2\le j\le h$, the recurrence denominator is at least $d$, and its two numerator coefficients are at most $A,v$. Induction therefore gives
\[
\mu_j\le\frac{A\Gamma^{j-1}+v\Gamma^{j-2}}d=\Gamma^j.
\]
The comparison with $\beta$ follows from
\begin{equation}\label{eq:beta}
(n-a+1)\beta^2=2(a-1)\beta+(a-1).
\end{equation}
Both right-hand coefficients increase strictly when $h<a$ is replaced by $h=a$, and the positive root increases with either coefficient. At $h=1$, the defining formula directly gives the stated first moment.
\end{proof}

\begin{corollary}[Scalar mixed bounds]\label{cor:mixedscalar}
For $h\ge1$,
\begin{equation}\label{eq:mixedgamma}
|\E P|\le 2^b n^{-a}F_n(a-h)\Gamma_{n,a,h}^{\,h}.
\end{equation}
For $2\le a\le n$ and every $h$,
\begin{equation}\label{eq:mixedgeneral}
|\E P|\le F_n(a-h)\,2^b n^{-a}\beta^h
\le\exp\left(\frac{a(a-1)}{2(n-a+1)}\right)2^bn^{-a}\beta^h.
\end{equation}
If $2\le a\le\sqrt n$, then
\begin{equation}\label{eq:mixedmain}
|\E P|\le\sqrt e\,2^b n^{-a}\beta^h
\le\sqrt e\,2^b n^{-a}\left(2\sqrt{a/n}\right)^h.
\end{equation}
The universal prefactor $\sqrt e$ in \eqref{eq:mixedmain} is optimal for this class.
\end{corollary}
\begin{proof}
The exact bound is $2^bn^{-a}F_n(a-h)\mu_h(n,a-h)$. Apply Lemma~\ref{lem:super}, with $\mu_0=1$ when $h=0$.

The elementary raw factorial estimate giving $\sqrt e$ is already recorded by Bordenave and Zhu \cite[Lemma 2.1, equation (2.3)]{BZ}. Here, for $0\le q\le a$, the inequality $-\log(1-x)\le x/(1-x)$ gives
\[
\log F_n(q)\le\frac{q(q-1)}{2(n-q+1)}
\le\frac{a(a-1)}{2(n-a+1)}.
\]
This is one global factor after the color merge. If $a^2\le n$, the last expression is at most $1/2$. Also $\sqrt{(a-1)/n}\le1/2$, so $\beta\le2\sqrt{a/n}$.

For sharpness, take $n=a^2$ and $a$ compatible raw cells in one color. Then $h=b=0$ and $n^a\E P=F_{a^2}(a)$, while
\[
\log F_{a^2}(a)=
\sum_{j=0}^{a-1}-\log(1-j/a^2)
=\frac{a(a-1)}{2a^2}+O(a^{-1})\longrightarrow\frac12.
\]
Thus no smaller universal prefactor works in \eqref{eq:mixedmain}.
\end{proof}

\subsection{Dependence on the product's composition}

For each fixed $h\ge1$, if $n\to\infty$ and $r/\sqrt n\to\rho\ge0$, then
\begin{equation}\label{eq:mixedfixedh}
n^{h/2}\mu_h(n,r)\longrightarrow
g_h(\rho)=h!\sum_{q=0}^{\lfloor h/2\rfloor}
\frac{\rho^{h-2q}}{(h-2q)!\,2^q q!}.
\end{equation}
Indeed, in the finite partition sum \eqref{eq:depletedpartitions}, singleton and pair weights are $r$ and $n+r$. After division by $n^{h/2}$ they have the displayed limits. Every block of size $\ell\ge3$ contributes a vanishing factor $O(n^{1-\ell/2})$. Finally, $(n-r)_h/n^h\to1$.

For example, take $n=a^2$, $r=a-h$ and fixed $h$ as $a\to\infty$. Then
\[
\frac{\mu_h(a^2,a-h)}{\beta^h}
\sim g_h(1)a^{-h/2},
\qquad
\Gamma_{a^2,a,h}\sim\frac{1+\sqrt{4h-3}}{2a}.
\]
This compares the full bounds, whose other factors coincide. The attaining family in Theorem~\ref{thm:mixed} realizes the improved order.

Such an improvement is not uniform when all singleton counts grow. For even $h=a$, pair partitions give
\[
\mu_a(n,0)\ge(a-1)!!\,n^{-a/2}
\ge\left(\frac{a}{en}\right)^{a/2}.
\]
The second inequality follows by summing the midpoint bound for $\log(2x)$ over $[0,a/2]$. Hence, along even $a\to\infty$ with $a^2\le n$, no fixed universal prefactor permits a singleton scale $o(\sqrt{a/n})$.

Finally, colors may have separate dimensions. If color $c$ has dimension $n_c\ge2$, centering $1/n_c$, and local counts $a_c,h_c,b_c$ with $a_c\le n_c$, independence gives
\begin{equation}\label{eq:mixed-dimensions}
|\E P|\le\prod_c
\left[
2^{b_c}\frac{n_c^{-h_c}\mathcal P_{h_c}(n_c,a_c-h_c)}
{(n_c)_{a_c}}
\right].
\end{equation}
There is no restriction on $\sum_c a_c$. If every $b_c=0$, compatible centered and raw choices attain all local bounds simultaneously, so \eqref{eq:mixed-dimensions} is again an exact maximum for the specified local counts.
